\section{Update Regimes in Lifelong Editing}

From a systems perspective, lifelong editing can be organized under at least two update regimes. The first is sequential updating, in which each new fact is incorporated as soon as it arrives. This regime is the most natural from a deployment viewpoint because it matches environments where updates are unpredictable and must take effect quickly. Its main weakness is that each new update is applied on top of the already edited model, which makes the procedure sensitive to interference and to the order in which edits arrive.

The second regime is batch updating, in which several edits are accumulated and optimized together. Batch editing is attractive because a joint update can satisfy several constraints at once rather than repeatedly perturbing the model in a one-by-one fashion. This often leads to better stability, especially when the editing method is designed to solve a joint optimization problem. However, batch updating also makes a stronger assumption about timing: it delays integration until several edits have been gathered, which is not always compatible with environments that require immediate response.

The distinction between sequential and batch updating is therefore not merely procedural. It changes what kind of editing method is favored. Methods optimized for joint consolidation may achieve stronger long-term factual integration, while methods optimized for immediate insertion may be more suitable for online use but less stable over long edit streams. In other words, the update regime is itself part of the problem definition rather than a minor implementation detail. This distinction forms part of the background needed to interpret the empirical results later in the thesis.

This observation is especially relevant for the current work because several strong editors appear most reliable when several edits can be considered together, while realistic deployment environments usually reveal edits one at a time or in small bursts. A useful lifelong editing system must therefore be evaluated not only by how well it rewrites facts, but also by how naturally its update regime matches the conditions under which those facts actually arrive and need to influence model behavior.
